% Changed: argument and layout integration; measured claims retained.
% Protected: existing citations, equations and labels unless explicitly documented.
% Meaning-Check: PRESERVED
% Risk-Flags: none
\section{Service Model and Design Objectives}
\label{sec:model}

\subsection{Structured Requests and Registered Dependencies}
A structured request $r=(s_1,\ldots,s_m)$ contains application segments. Each segment carries an identifier, original content, a stability type, and a registered content identity. We retain three types: \textsc{Stable}, \textsc{Mutable}, and \textsc{Query}. Stable content belongs to an immutable template version; mutable content includes observations and changing state; query content carries the current task. A query may recur within an authorized directory while remaining outside the promoted stable prefix.

The operator manifest defines segment boundaries, separators, stability, dependencies, and the sharing domain. Unregistered tenant content remains private. A segment dependency graph $G_r=(V,E,\lambda)$ associates each edge with a data, causal, or positional ordering constraint. A layout $\pi$ is legal when it is a topological order of $G_r$:
\begin{equation}
 (s_i,s_j)\in E\quad\Longrightarrow\quad
 \operatorname{pos}_{\pi}(s_i)<\operatorname{pos}_{\pi}(s_j).
\end{equation}

All cache identities and prefix boundaries are computed from the fully rendered prompt $\mathbf{t}_{\pi}$, including the model's chat template. Per-segment token counts guide construction, but exact boundaries follow retokenization; concatenation can change tokenization at segment edges. Let $P(\pi)$ denote the token boundary of the leading shareable stable content. CALC provides a bounded candidate set $\Pi_L(G_r)$ containing raw order L0 and registered compiled levels. Behavioral qualification selects among these legal candidates.

\subsection{Templates, Qualifications, and Cache Identity}
A service template
\begin{equation}
 \Theta=(V,\mathbf S,\mathcal A,\mathcal P,\mathcal M,F)
\end{equation}
binds a version $V$, stable content $\mathbf S$, an authorized tenant set $\mathcal A$, a frozen policy $\mathcal P$, its manifest $\mathcal M$, and declared fan-out $F$. The input may interleave stable and private segments; stable-first order is a compilation result.

ContractKV uses two related identities. A \emph{behavioral qualification} determines whether a request may use a particular compiled layout. A \emph{stable-prefix identity} identifies the corresponding reusable KV prefix within its authorized template namespace. Several admitted query scopes can map to the same stable-prefix identity.

Execution identity binds the model, tokenizer/chat template, decoding/stopping settings, policy and runtime. Section~\ref{sec:template} defines the complete template and lifecycle binding.

\subsection{Publication and Service Objectives}
ContractKV aims to turn admissible structural reuse into lower complete service cost while preserving the following publication properties.

\noindent\textbf{P1: Evidence-gated publication.} Candidate compiled state becomes shared only after a completed verifier pass and a valid publication binding.

\noindent\textbf{P2: Stable-prefix confinement.} PSTP publishes complete blocks within the stable, certified, exact-match boundary. Mutable and query suffixes are excluded from this promotion.

\noindent\textbf{P3: Atomic failure handling.} A failed, incomplete, expired, or conflicting transaction contributes no newly promoted candidate blocks. Raw-order execution remains the service fallback.

\noindent\textbf{P4: Authorization-bound reuse.} The deployment control path validates template identity and authorization before admitting reuse, and separates cache namespaces across incompatible sharing domains.

These properties concern admission and publication. We assess behavior with paired free-running outputs and task scores, and efficiency with native hits and complete service wall time.
